% Source: Wald GR, physical PDF pp. 337--343 (printed pp. 324--330).
% Coverage: Section 12.4, Energy Extraction from Black Holes; equations (12.4.1)--(12.4.24).
% Figures/tables/footnotes: Figures 12.7--12.8; no tables or footnotes.
% Status: source-reviewed and compile-clean.
% Uncertainties: none.

\subsection{Energy Extraction from Black Holes}\label{sec:12.4}
\setcounter{theorem}{0}

By definition, a black hole is a \lq\lq region of no escape.\rq\rq\ No material body or
light ray ever can be extracted from a black hole. Therefore, it came as a great
surprise when Penrose (1969) noted that energy can be extracted from a black hole
with an ergosphere. The mechanism proposed by Penrose can be understood as
follows. The Killing field \(\xi^a\) which becomes a time translation asymptotically
at infinity is spacelike in the ergosphere. Thus, for a test particle of 4-momentum
\(p^a=mu^a\), the energy
\begin{equation}
 E=-p^a\xi_a
 \tag{12.4.1}\label{eq:12.4.1}
\end{equation}
need not be positive in the ergosphere. Therefore, by making a black hole absorb a
particle with negative total energy, we can extract energy from a black hole! To see
this in more detail, suppose we start in our laboratory far from the black hole by
throwing a particle toward the black hole. If we denote the 4-momentum of this
particle by \(p_0^a\), its total energy measured in the laboratory will be
\begin{equation}
 E_0=-p_0^a\xi_a .
 \tag{12.4.2}\label{eq:12.4.2}
\end{equation}
As it falls freely toward the black hole, \(E_0\) will remain constant. Suppose,
when the particle enters the ergosphere, we arrange---e.g., by means of explosives
and a timing device---to have it break up into two fragments as illustrated in
Figure~12.7. By local conservation of energy-momentum, we have
\begin{equation}
 p_0^a=p_1^a+p_2^a ,
 \tag{12.4.3}\label{eq:12.4.3}
\end{equation}
where \(p_1^a\) and \(p_2^a\) are the 4-momenta of the two fragments. Contracting
equation~\eqref{eq:12.4.3} with \(\xi_a\), we obtain
\begin{equation}
 E_0=E_1+E_2 .
 \tag{12.4.4}\label{eq:12.4.4}
\end{equation}
However, inside the ergosphere, we can arrange the breakup so that one of the
fragments has negative total energy,
\begin{equation}
 E_1<0 .
 \tag{12.4.5}\label{eq:12.4.5}
\end{equation}
Therefore, if the other fragment returns to our laboratory in free (i.e., geodesic)
motion, it will have an energy \(E_2\) which is greater than the initial energy \(E_0\).

In the case of a Kerr black hole of mass \(M\) with \(a\ne0\), one may explicitly
verify that the breakup process can be done so that the second fragment does,
indeed, escape to infinity. One also may verify that the negative energy fragment
always falls into the black hole. Thus, at the end of the process, one has energy
\(E_0+|E_1|\) in the laboratory, and the mass of the black hole must be reduced to
\(M-|E_1|\). Thus, the energy \(|E_1|\) has been extracted from the black hole!

How much energy can be extracted from a Kerr black hole in this manner? As we
shall see below shortly, the energy extraction process is self-limiting because the
negative energy particles which enter the black hole also carry negative angular
momentum, i.e., angular momentum opposite that of the black hole. As a result, the
angular momentum \(J=Ma\) of the black hole will be reduced to zero while \(M\) is
still finite. However, when \(J=0\) the ergosphere no longer is present, and no
further energy extraction can occur.

% Source: Wald GR, physical PDF p. 338
%         (printed p. 325).
% Coverage: The Penrose energy-extraction process.
% Figures/tables/footnotes: Figure 12.7; no tables or footnotes.
% Status: source-reviewed and compile-clean.
% Uncertainties: none.

\begingroup
\renewcommand{\thefigure}{12.7}
\begin{figure}[H]
  \centering
  \begin{tikzpicture}[x=.78cm,y=.78cm,>=Latex,line cap=round,line join=round,font=\small]
    \draw[thick] (0,0) circle [radius=.72];
    \draw[thick] (0,0) circle [radius=1.32];
    \node[align=center] at (0,0) {Black\\Hole};
    \draw[->,thick] (-4.0,.25) .. controls (-2.55,-.62) and (-1.52,-.94) .. (-.35,-.82);
    \draw[->,thick] (.00,-.95) .. controls (1.15,-1.42) and (2.55,-1.48) .. (3.75,-.85);
    \draw[->,thick] (-.18,-.86) -- (.10,-.52);
    \node at (-2.40,-1.03) {\(E_0\)};
    \node at (2.68,-1.70) {\(E_2\)};
    \node[right] at (.34,-.76) {\(E_1\)};
    \draw[->] (1.85,1.10) .. controls (1.52,.65) and (1.24,.38) .. (.95,.26);
    \node[right] at (1.80,1.17) {Ergosphere};
    \draw[->,thick] (.55,.94) arc[start angle=35,end angle=145,x radius=.62,y radius=.22];
  \end{tikzpicture}
  \caption{A diagram illustrating the Penrose process for extracting energy from a
  Kerr black hole.}
  \label{fig:12.7}
\end{figure}
\endgroup


To see this limit on energy extraction in detail, we use the fact that the Killing field
\(\chi^a\), defined by equation~(12.3.20), is future directed null on the horizon.
Hence for any particle which enters the black hole (which includes all negative
energy particles), we have
\begin{equation}
 0>p^a\chi_a=p^a(\xi_a+\Omega_H\psi_a)=-E+\Omega_HL ,
 \tag{12.4.6}\label{eq:12.4.6}
\end{equation}
where \(L=p^a\psi_a\) and \(\Omega_H\) was defined by equation~(12.3.19). Thus,
we find that
\begin{equation}
 L<E/\Omega_H ,
 \tag{12.4.7}\label{eq:12.4.7}
\end{equation}
which verifies the above statement that a negative-energy particle entering the black
hole carries negative angular momentum. After the black hole \lq\lq swallows\rq\rq\ a
particle, it should settle back down to a Kerr solution with parameters modified by
\(\delta M=E\), \(\delta J=L\). Thus, from equation~\eqref{eq:12.4.7}, the change
in black hole parameters is restricted by
\begin{equation}
 \delta J<\delta M/\Omega_H
 \tag{12.4.8}\label{eq:12.4.8}
\end{equation}
which can be rewritten as (Christodoulou 1970)
\begin{equation}
 \delta M_{\mathrm{irr}}>0 ,
 \tag{12.4.9}\label{eq:12.4.9}
\end{equation}
where the \emph{irreducible mass}, \(M_{\mathrm{irr}}\), is defined by
\begin{equation}
 M_{\mathrm{irr}}^2=\frac12\left[M^2+(M^4-J^2)^{1/2}\right] .
 \tag{12.4.10}\label{eq:12.4.10}
\end{equation}
Inverting equation~\eqref{eq:12.4.10}, we find
\begin{equation}
 M^2=M_{\mathrm{irr}}^2+\frac14\frac{J^2}{M_{\mathrm{irr}}^2}
 \geq M_{\mathrm{irr}}^2 .
 \tag{12.4.11}\label{eq:12.4.11}
\end{equation}
Thus, the mass of a black hole cannot be reduced below the initial value of
\(M_{\mathrm{irr}}\) via the Penrose process. If we start with a Kerr black hole of
mass \(M_0\) and angular momentum \(J_0\), then even assuming that equality is
achieved in equation~\eqref{eq:12.4.7}, by the time the energy
\(M_0-M_{\mathrm{irr}}(M_0,J_0)\) has been extracted, the angular momentum of the
black hole will have been reduced to zero. Since one can come arbitrarily close to
achieving equality in equation~\eqref{eq:12.4.7}, we can come arbitrarily close to
extracting energy \(M_0-M_{\mathrm{irr}}\) from the black hole via the Penrose
process. We may interpret \(M_0-M_{\mathrm{irr}}\) as the rotational energy of the
black hole. For a maximally rotating black hole, \(J_0=M_0^2\), it represents
\((1-1/\sqrt2)\approx29\%\) of the mass-energy of the black hole.

The universal nature of the limit on energy extraction implied by
equation~\eqref{eq:12.4.11}---obtained above only for the specific process proposed
by Penrose---can be seen from area theorem~12.2.6. The area of the event horizon
of a Kerr black hole is given by
\begin{equation}
\begin{aligned}
 A&=\int_{r=r_+}\sqrt{g_{\theta\theta}g_{\phi\phi}}\,\dd\theta\,\dd\phi\\
  &=\int(r_+^2+a^2)\sin\theta\,\dd\theta\,\dd\phi\\
  &=4\pi(r_+^2+a^2)\\
  &=16\pi M_{\mathrm{irr}}^2 .
\end{aligned}
\tag{12.4.12}\label{eq:12.4.12}
\end{equation}
Thus, from the area theorem we obtain the completely general result that
\(M_{\mathrm{irr}}\) can never decrease, from which follows the above energy
extraction limit for all possible processes. The consistency of the area theorem,
proven by general, abstract arguments, with the result~\eqref{eq:12.4.11} obtained
by the calculation of a specific process for a Kerr black hole, lends further aesthetic
support for the validity of the first cosmic censor conjecture, which is used in the
above argument when we assume that the infalling particle merely changes the black
hole parameters and does not convert the black hole to a naked singularity.

It should be noted that the area theorem also can be used to obtain an interesting
upper limit on the energy radiated away in the form of gravitational waves when two
black holes collide (Hawking 1971). Consider, for example, initial data representing
two widely separated Schwarzschild black holes initially \lq\lq at rest,\rq\rq\ with
masses \(M_1\) and \(M_2\). (See the paragraph below equation~[10.2.39] of
Section~10.2 for a brief discussion of how to construct such data explicitly.)
Presumably, the dynamical evolution of these data will yield a spacetime where the
two black holes fall toward each other, coalesce, and \lq\lq settle down\rq\rq\ to a
single Schwarzschild black hole of mass \(M\). The total energy radiated away in
this process is
\begin{equation}
 E_{\mathrm{rad}}=M_1+M_2-M .
 \tag{12.4.13}\label{eq:12.4.13}
\end{equation}
An upper limit on \(E_{\mathrm{rad}}\) can be obtained by noting that the initial black
hole area is
\begin{equation}
 A_i=A_1+A_2=16\pi(M_1^2+M_2^2) .
 \tag{12.4.14}\label{eq:12.4.14}
\end{equation}
The final area is
\begin{equation}
 A_f=16\pi M^2 ;
 \tag{12.4.15}\label{eq:12.4.15}
\end{equation}
and, by the area theorem, we have
\begin{equation}
 A_f\geq A_i .
 \tag{12.4.16}\label{eq:12.4.16}
\end{equation}
Putting together equations~\eqref{eq:12.4.13}--\eqref{eq:12.4.16}, we obtain
\begin{equation}
 E_{\mathrm{rad}}\leq M_1+M_2-(M_1^2+M_2^2)^{1/2} .
 \tag{12.4.17}\label{eq:12.4.17}
\end{equation}
For the case \(M_1=M_2\), this implies that at most
\((1-1/\sqrt2)\approx29\%\) of the original mass can be radiated away. Numerical
calculations (Smarr 1979) indicate that far less energy than this upper limit actually
will be radiated in this process.

Although the Penrose process is of great importance for demonstrating that, in
principle, the maximum amount of energy permitted by the area theorem can be
extracted from a rotating black hole, the process requires a precisely timed breakup
of the incident particle at relativistic velocities and it is not a practical energy
extraction method (Bardeen, Press, and Teukolsky 1972; Wald 1974c). Interestingly
there is a wave analog of the Penrose process, known as \emph{superradiant
scattering} (Misner, unpublished; Zel'dovich 1972; Starobinskii 1973), which allows
energy to be extracted from a black hole in a relatively simple manner. If a scalar,
electromagnetic, or gravitational wave is incident upon a black hole, part of the wave
(the \lq\lq transmitted wave\rq\rq ) will be absorbed by the black hole and part of the
wave (the \lq\lq reflected wave\rq\rq ) will escape back to infinity. Normally the
transmitted wave will carry positive energy into the black hole, and the reflected wave
will have less energy than the incident wave. However, for a wave of the form
\(\phi=\operatorname{Re}[\phi_0(r,\theta)e^{-\ii\omega t}e^{\ii m\phi}]\) with
\begin{equation}
 0<\omega<m\Omega_H
 \tag{12.4.18}\label{eq:12.4.18}
\end{equation}
the transmitted wave will carry negative energy into the black hole (analogous to the
negative energy fragment in the Penrose process for particles) and the reflected wave
will return to infinity with greater amplitude and energy than the incident wave. This
is demonstrated most easily for the case of scalar waves. By contracting the stress
tensor, equation~(4.3.10), of a Klein--Gordon scalar field \(\phi\) with the timelike
Killing field \(\xi^a\) of the Kerr background, we obtain an \lq\lq energy current,\rq\rq
\begin{equation}
 f^a=-T^a{}_{b}\xi^b
 \tag{12.4.19}\label{eq:12.4.19}
\end{equation}
which is conserved since \(\nabla_af^a=-(\nabla_aT^a{}_{b})\xi^b-
T^{ab}\nabla_a\xi_b=0\). Hence, if we integrate \(\nabla_af^a\) over the region
\(K\) of spacetime shown in Figure~12.8, we find by Gauss's law that the difference
between the incoming and outgoing energies (i.e., the integrated flux of \(f^a\) over
the \lq\lq large sphere\rq\rq ) equals the integrated flux of \(f^a\) on the horizon.
However, on the horizon the time averaged flux is given by
\begin{equation}
\begin{aligned}
 \langle f^an_a\rangle&=-(T_{ab}\xi^b)\chi^a=(T_{ab}\chi^a\xi^b)\\
 &=\langle(\chi^a\nabla_a\phi)(\xi^b\nabla_b\phi)\rangle\\
 &=\frac12\omega(\omega-m\Omega_H)|\phi_0|^2 .
\end{aligned}
\tag{12.4.20}\label{eq:12.4.20}
\end{equation}

% Source: Wald GR, physical PDF p. 341
%         (printed p. 328).
% Coverage: Spacetime region used in the superradiance flux calculation.
% Figures/tables/footnotes: Figure 12.8; no tables or footnotes.
% Status: source-reviewed and compile-clean.
% Uncertainties: none.

\begingroup
\renewcommand{\thefigure}{12.8}
\begin{figure}[H]
  \centering
  \begin{tikzpicture}[x=.85cm,y=.85cm,>=Latex,line cap=round,line join=round,font=\small]
    \coordinate (Hone) at (-2.05,-.95);
    \coordinate (Htwo) at (-.58,.65);
    \coordinate (Sone) at (2.10,-.74);
    \coordinate (Stwo) at (2.38,.54);
    \path[fill=black!7] (Hone)
      .. controls (-.65,-1.04) and (.92,-1.16) .. (Sone)
      -- (Stwo)
      .. controls (1.15,.28) and (.20,.38) .. (Htwo) -- cycle;
    \draw[very thick] (-2.72,-1.68) -- (.28,1.64);
    \draw[thick] (Hone) .. controls (-.65,-1.04) and (.92,-1.16) .. (Sone);
    \draw[thick] (Htwo) .. controls (.20,.38) and (1.15,.28) .. (Stwo);
    \draw[thick] (Sone) -- (Stwo);
    \node[left] at (-1.92,.18) {Horizon};
    \node[right] at (2.30,-.10) {\(S\)};
    \node at (.20,-1.19) {\(\Sigma_1\)};
    \node at (.52,.70) {\(\Sigma_2\)};
    \node at (.55,-.35) {\(K\)};
    % Outward normals on all four pieces of the boundary.
    \draw[->] (-1.53,-.38) -- (-1.84,-.72) node[midway,left] {\(n^a\)};
    \draw[->] (1.04,.32) -- (.98,.00) node[midway,right] {\(n^a\)};
    \draw[->] (.10,-1.04) -- (.14,-.68) node[midway,right] {\(n^a\)};
    \draw[->] (2.26,.16) -- (2.66,.22);
    \node[above] at (2.58,.30) {\(n^a\)};
  \end{tikzpicture}
  \caption{A spacetime diagram showing the region \(K\) over which
  \(\nabla_a f^a=0\) is integrated to derive the conclusion in the text concerning
  superradiant scattering. Here the spacelike hypersurface \(\Sigma_2\) is a
  \lq\lq time translate\rq\rq\ of \(\Sigma_1\) by \(\Delta t\), and the timelike
  hypersurface \(S\) represents a \lq\lq large sphere\rq\rq\ at infinity. The integral
  of \(f^a n_a\) over \(S\) represents the net energy flow out of \(K\) to infinity
  (i.e., the outgoing minus incoming energy) during the time \(\Delta t\), whereas
  the integral of \(f^a n_a\) over the horizon represents the net energy flow into the
  black hole. By Gauss's law (see appendix B) the integral of \(f^a n_a\) over the
  entire boundary vanishes. (The appropriate directions of \(n^a\) on each part of
  the boundary are shown.) However, for a wave with time dependence \(e^{-\ii\omega t}\)
  the integrals over \(\Sigma_1\) and \(\Sigma_2\) cancel by time translation symmetry.
  Hence, the integral of \(f^a n_a\) over \(S\) equals minus the integral of
  \(f^a n_a\) over the horizon.}
  \label{fig:12.8}
\end{figure}
\endgroup


where \(n^a=-\chi^a\) (with \(\chi^a\) given by equation~[12.3.20]) is the
appropriately directed normal to the horizon. (The term in \(T_{ab}\) proportional to
\(g_{ab}\) does not contribute since \(\chi^a\xi_a=0\) on the horizon.) Thus, in the
frequency range of equation~\eqref{eq:12.4.18} the energy flux through the horizon
is negative, and hence superradiance is obtained. This conclusion also can be derived
by considering the \lq\lq particle number current,\rq\rq
\begin{equation}
 j^a=-\ii(\bar\phi\nabla^a\phi-\phi\nabla^a\bar\phi)
 \tag{12.4.21}\label{eq:12.4.21}
\end{equation}
for the complex field \(\phi_0e^{-\ii\omega t}e^{\ii m\phi}\). This current is
conserved by virtue of the Klein--Gordon equation. Again, \(-j^a\chi_a\) is negative
in the range~\eqref{eq:12.4.18}, thus establishing superradiance.

The superradiance of electromagnetic waves in the range~\eqref{eq:12.4.18} can be
proven similarly by consideration of the energy current, \(-T_{ab}\xi^b\), or a
(non-gauge invariant) \lq\lq particle number current\rq\rq\ analogous to
equation~\eqref{eq:12.4.21}, although the demonstration of the negative integrated
flux of the energy current through the horizon is less straightforward (see problem
5). In the gravitational case, we also can form a conserved \lq\lq effective energy
current,\rq\rq\ \((1/8\pi)G^{(2)}_{ab}\xi^b\), where \(G^{(2)}_{ab}\) is the second
order Einstein tensor (see Section~4.4b for the expression for \(G^{(2)}_{ab}\) for
perturbations off flat spacetime), and a conserved \lq\lq particle number current.\rq\rq
The formula for the energy current is rather complicated, and neither current is gauge
invariant. However, one can avoid dealing with them by working, instead, with the
decoupled Teukolsky equation for the variables \(\Psi_0\) or \(\Psi_4\) defined in
Section~12.3. By doing so, superradiance in the range~\eqref{eq:12.4.18} for
gravitational waves can be established in a relatively straightforward manner
(Teukolsky and Press 1974). Superradiance of electromagnetic waves also can be
proven from the Teukolsky equation for \(\Phi_0\) or \(\Phi_2\). Numerical
computations (Teukolsky and Press 1974) have shown that the largest superradiance
effects occur for gravitational waves, where an amplification factor of up to \(1.38\)
can be achieved for a Kerr black hole with \(a=M\).

The necessity of superradiance in the regime~\eqref{eq:12.4.18} also can be seen
directly from the area theorem (Bekenstein 1973a). For a wave with frequency
\(\omega\) and azimuthal number \(m\) in a stationary axisymmetric background, the
ratio of angular momentum flux to energy flux at infinity is
\begin{equation}
 \mathcal L/\mathcal E=m/\omega .
 \tag{12.4.22}\label{eq:12.4.22}
\end{equation}
Thus, by conservation of energy and angular momentum, when such a wave is
incident upon a black hole, the change in energy, \(\delta M\), and angular momentum,
\(\delta J\), of the black hole must be related by
\begin{equation}
 \delta J/\delta M=m/\omega .
 \tag{12.4.23}\label{eq:12.4.23}
\end{equation}
However, in the range~\eqref{eq:12.4.18}, we have
\begin{equation}
 \delta J/\delta M>1/\Omega_H .
 \tag{12.4.24}\label{eq:12.4.24}
\end{equation}
If \(\delta M>0\), this would violate equation~\eqref{eq:12.4.8}, which, by
equations~\eqref{eq:12.4.9} and~\eqref{eq:12.4.12}, would violate the area theorem.
Hence, we must have \(\delta M<0\)---i.e., superradiance must occur---when
\(0<\omega<m\Omega_H\).

Interestingly, fermion fields do \emph{not} display superradiance (Unruh 1973;
G\"uven 1977). The \lq\lq particle number current,\rq\rq\ \(j^a\), associated with
neutrino or Dirac fields (see chapter 13) is manifestly a timelike or null vector, and
hence \(j_an^a\) is manifestly nonnegative on the horizon for all waves including
those in the range~\eqref{eq:12.4.18}. Thus, the reflected wave never has larger
amplitude than the incident wave. In this case the argument given in the previous
paragraph is inapplicable because the stress tensor of these fields fails to satisfy the
weak energy condition, so the area theorem does not hold.

The behavior of both boson and fermion fields incident upon a Kerr black hole with
\(a\ne0\) is in very close mathematical analogy to a well studied effect in
nongravitational physics known as the Klein \lq\lq paradox.\rq\rq\ If a Klein--Gordon
field of charge \(Q\) and mass \(M\) in one spatial dimension is incident upon an
electrostatic potential, \(V\), such that \(V\to0\) as \(x\to-\infty\) but
\(V\to V_0>2M/Q\) for \(x\to+\infty\), then when \(\omega+M<QV_0\) the reflected
wave has greater amplitude and energy than the incident wave. For a Dirac field, the
reflected wave will be smaller than the incident wave---again the Dirac \lq\lq particle
number current\rq\rq\ is always timelike---but the transmitted wave still has a
negative \lq\lq kinetic energy\rq\rq\ (Klein 1929). In quantum field theory, the
interpretation of the Klein paradox is that in both the boson and fermion cases
particle-antiparticle pairs are spontaneously created in the strong electrostatic field
associated with the potential, \(V\). When incoming particles also are present,
stimulated emission occurs in the boson case, and in the classical limit this gives rise
to the larger \lq\lq reflected wave\rq\rq\ obtained in the classical analysis. The close
analogy between the Klein paradox and the scattering of waves by a Kerr black hole
suggests that spontaneous particle creation should occur near a Kerr black hole.
Indeed, this is the case, but we shall postpone our discussion of this phenomenon
until chapter 14.
